\titledquestion{Machine Learning \& Neural Networks}[8]
\begin{parts}


    \part[4] Adam Optimizer\newline
    Recall the standard Stochastic Gradient Descent update rule:
    \alns{
        \btheta_{t+1} &\gets \btheta_t - \alpha \nabla_{\btheta_t} J_{\text{minibatch}}(\btheta_t)
    }
    where $t+1$ is the current timestep, $\btheta$ is a vector containing all of the model parameters, ($\btheta_t$ is the model parameter at time step $t$, and $\btheta_{t+1}$ is the model parameter at time step $t+1$), $J$ is the loss function, $\nabla_{\btheta} J_{\text{minibatch}}(\btheta)$ is the gradient of the loss function with respect to the parameters on a minibatch of data, and $\alpha$ is the learning rate.
    Adam Optimization\footnote{Kingma and Ba, 2015, \url{https://arxiv.org/pdf/1412.6980.pdf}} uses a more sophisticated update rule with two additional steps.\footnote{The actual Adam update uses a few additional tricks that are less important, but we won't worry about them here. If you want to learn more about it, you can take a look at: \url{http://cs231n.github.io/neural-networks-3/\#sgd}}

    \begin{subparts}

        \subpart[2]First, Adam uses a trick called {\it momentum} by keeping track of $\bm$, a rolling average of the gradients:
        \alns{
            \bm_{t+1} &\gets \beta_1\bm_{t} + (1 - \beta_1)\nabla_{\btheta_t} J_{\text{minibatch}}(\btheta_t) \\
            \btheta_{t+1} &\gets \btheta_t - \alpha \bm_{t+1}
        }
        where $\beta_1$ is a hyperparameter between 0 and 1 (often set to  0.9). Briefly explain in 2--4 sentences (you don't need to prove mathematically, just give an intuition) how using $\bm$ stops the updates from varying as much and why this low variance may be helpful to learning, overall.\newline


        \subpart[2] Adam extends the idea of {\it momentum} with the trick of {\it adaptive learning rates} by keeping track of  $\bv$, a rolling average of the magnitudes of the gradients:
        \alns{
            \bm_{t+1} &\gets \beta_1\bm_{t} + (1 - \beta_1)\nabla_{\btheta_t} J_{\text{minibatch}}(\btheta_t) \\
            \bv_{t+1} &\gets \beta_2\bv_{t} + (1 - \beta_2) (\nabla_{\btheta_t} J_{\text{minibatch}}(\btheta_t) \odot \nabla_{\btheta_t} J_{\text{minibatch}}(\btheta_t)) \\
            \btheta_{t+1} &\gets \btheta_t - \alpha \bm_{t+1} / \sqrt{\bv_{t+1}}
        }
        where $\odot$ and $/$ denote elementwise multiplication and division (so $\bz \odot \bz$ is elementwise squaring) and $\beta_2$ is a hyperparameter between 0 and 1 (often set to  0.99). Since Adam divides the update by $\sqrt{\bv}$, which of the model parameters will get larger updates?  Why might this help with learning?



    \end{subparts}
    \ifans{
        \begin{enumerate}[label=\roman*.]
            \item Since the hyperparameter $\beta_1$ often set to 0.9, $\bm$ receive most of its value from the last step and a little from the gradient. Therefore the update has a sense of inertia that avoids the magnitude and direction change too much. The low variance can make a faster convergence to the nearest optimum.

                  $\downarrow\downarrow\downarrow$maybe a better answer$\downarrow\downarrow\downarrow$

                  Instead of updating model parameters based solely on the current gradient, $momentum$ smooths out the updates by accumulating an exponentially decaying average of past gradients. This approach prevents sudden shifts in the parameters due to noisy gradients and reduces oscillations when gradients fluctuate in direction. Additionally, momentum helps the parameter update maintain speed, which prevents learning from getting stuck in plateaus or local minima.

            \item The parameters which have smaller gradient will get larger updates. Instead of applying the same learning rate to all parameters, the scaling operation $\slash \sqrt{\bv}$ adapts the learning rate for each parameter. This allows parameters with infrequent or small gradients to be updated with larger step sizes, leading to faster convergence. Conversely, parameters with larger gradient magnitudes receive smaller updates, preventing overshooting and stabilizing the learning process. Overall, $adaptive learning rates$ help control parameter-wise updates, enhancing the model's generalizability.
        \end{enumerate}

    }

    \part[4]
    Dropout\footnote{Srivastava et al., 2014, \url{https://www.cs.toronto.edu/~hinton/absps/JMLRdropout.pdf}} is a regularization technique. During training, dropout randomly sets units in the hidden layer $\bh$ to zero with probability $p_{\text{drop}}$ (dropping different units each minibatch), and then multiplies $\bh$ by a constant $\gamma$. We can write this as:
    \alns{
        \bh_{\text{drop}} = \gamma \bd \odot \bh
    }
    where $\bd \in \{0, 1\}^{D_h}$ ($D_h$ is the size of $\bh$)
    is a mask vector where each entry is 0 with probability $p_{\text{drop}}$ and 1 with probability $(1 - p_{\text{drop}})$. $\gamma$ is chosen such that the expected value of $\bh_{\text{drop}}$ is $\bh$:
    \alns{
    \mathbb{E}_{p_{\text{drop}}}[\bh_\text{drop}]_i = h_i \text{\phantom{aaaa}}
    }
    for all $i \in \{1,\dots,D_h\}$.
    \begin{subparts}
        \subpart[2]
        What must $\gamma$ equal in terms of $p_{\text{drop}}$? Briefly justify your answer or show your math derivation using the equations given above.


        \subpart[2] Why should dropout be applied during training? Why should dropout \textbf{NOT} be applied during evaluation? \textbf{Hint:} it may help to look at the dropout paper linked. \newline


    \end{subparts}

    \ifans{
        \begin{enumerate}[label=\roman*.]
            \item
                  \begin{gather*}
                      \E_{p_{\text{drop}}}[\bh_\text{drop}]_i=h_i\\
                      \E_{p_{\text{drop}}}(\gamma d_i h_i)=h_i\\
                      \gamma h_i \E_{p_{\text{drop}}}(d_i)=h_i\\
                      \gamma (1-p_{\text{drop}})=1\\
                      \gamma=\frac{1}{1-p_{\text{drop}}}
                  \end{gather*}

            \item During training, randomly masking a portion of neurons helps alleviate a phenomenon, called \textbf{co-adaptation}, introduced in the paper. Concretely speaking, dropout prevents each neuron from heavily depending on a specific set of other neurons to correct the mistake, which can lead to better generalizability.
            
            For evaluation (or inference), dropout is disabled so all neurons contribute to the prediction, ensuring stable and consistent outputs. To align expected outputs at training and test time, one approach (called \textbf{inverted dropout}) scales activations during training by dividing the keep probability, while keeping test activations unchanged. Alternatively, training remains unchanged, and activations are scaled down during testing.
        \end{enumerate}
    }


\end{parts}
